% !TeX root = ../main.tex
% Appendix B1: Quantum Copula Bounds and the CHSH Frontier
% Batch #4.9c — Notation hygiene

\section{Quantum Copula Bounds and the CHSH Frontier}\label{app:copula-chsh}

\subsection{Classical Copula Bounds}

A classical copula $C : [0, 1]^2 \to [0, 1]$ is a bivariate distribution function
with uniform marginals, constrained by the Fr\'echet--Hoeffding bounds:
\begin{equation*}
    \max(u + v - 1, 0) \leq C(u, v) \leq \min(u, v).
\end{equation*}

In the context of financial contagion, copulas are used to model the joint
distribution of asset returns or default events. The classical bounds imply that
the correlation between financial variables is limited by their marginal
distributions.

\subsection{Quantum Violation of Classical Copula Bounds}

In the quantum framework, the joint distribution of two non-commuting observables
is not a classical copula but a \emph{quantum copula}, which can violate the
Fr\'echet--Hoeffding bounds. This violation is a consequence of the
non-commutative structure of the observable algebra and is analogous to the
violation of Bell inequalities in quantum mechanics
\citep{clauser1969proposed,fine1982hidden}.

\begin{theorem}[Quantum copula bound]\label{thm:quantum-copula}
For two non-commuting market observables $\hat{A}, \hat{B} \in \mathcal{A}$ with
$[\hat{A}, \hat{B}] \neq 0$, the quantum copula $C_Q$ satisfies
\begin{equation*}
    C_Q(u, v) \leq \min(u, v) + \frac{\hbar_{\mathrm{econ}}}{2}
    \cdot \|\hat{A}\| \cdot \|\hat{B}\|,
\end{equation*}
where $\|\cdot\|$ denotes the operator norm. The classical Fr\'echet--Hoeffding
upper bound is violated when $\hbar_{\mathrm{econ}} > 0$.
\end{theorem}

This result has direct implications for financial contagion modeling. The joint
probability of extreme co-movements implied by equity option prices can exceed
the classical copula bound when the market is subject to microstructural
uncertainty. This provides a theoretical explanation for the ``correlation
skew'' observed in equity index options, where the implied correlation exceeds
the bounds of classical copula models.

\subsection{The CHSH Frontier in Financial Markets}

The Clauser--Horne--Shimony--Holt (CHSH) inequality
\citep{clauser1969proposed} provides a quantitative test of local hidden variable
theories. In the financial context, the CHSH inequality tests whether the joint
distribution of two market observables can be explained by a classical latent
factor model.

Define the CHSH parameter for four market observables
$\hat{A}_1, \hat{A}_2, \hat{B}_1, \hat{B}_2$ as
\begin{equation*}
    S = \langle \hat{A}_1 \hat{B}_1 \rangle + \langle \hat{A}_1 \hat{B}_2 \rangle
    + \langle \hat{A}_2 \hat{B}_1 \rangle - \langle \hat{A}_2 \hat{B}_2 \rangle.
\end{equation*}

In a classical latent factor model, $|S| \leq 2$ (the CHSH bound). In the quantum
framework, the Tsirelson bound \citep{cirelson1980quantum} applies:
$|S| \leq 2\sqrt{2}$.

\begin{proposition}[Financial CHSH violation]\label{prop:chsh-violation}
For appropriately chosen equity option observables of financial institutions, the
CHSH parameter $S$ can exceed the classical bound of $2$ when
$\hbar_{\mathrm{econ}} > 0$, providing empirical evidence for the non-commutative
structure of financial markets.
\end{proposition}

The empirical test of this proposition is a direction for future research. A
finding that $S > 2$ in equity markets would provide direct evidence that the
joint distribution of financial asset returns cannot be fully captured by
classical latent factor models, supporting the quantum framework's prediction of
non-commutative market microstructure.
